% Corrigé de l'exercice 10_PompePalette.
% Version robuste pour MiKTeX : aucun environnement listings ou verbatim.
\begin{corrigebox}[Corrigé complété]

\CorrigeQuestion{1}
Le graphe des liaisons comporte le bâti \textbf{0}, le rotor \textbf{1} et la
palette \textbf{2}. Le rotor est en liaison pivot avec le bâti, la palette est
en liaison glissière avec le rotor et son extrémité reste en contact avec le
stator.
\begin{marginfigure}
\includegraphics[width=\linewidth]{10_01_c}
\end{marginfigure}

\CorrigeQuestion{2}
La fermeture géométrique s'écrit
\[
\overrightarrow{OA}+\overrightarrow{AB}+\overrightarrow{BO}=\vec 0.
\]
Avec $\overrightarrow{AO}=e\,\vec i_0$,
$\overrightarrow{AB}=\lambda\,\vec i_1$ et
$\overrightarrow{OB}=R\,\vec u$, on obtient, après projection,
\[
\left\{
\begin{aligned}
-e+\lambda\cos\theta&=R\cos\varphi,\\
\lambda\sin\theta&=R\sin\varphi.
\end{aligned}
\right.
\]
En élevant au carré puis en additionnant les deux relations :
\[
R^2=e^2-2e\lambda\cos\theta+\lambda^2.
\]
La solution correspondant au contact avec le stator est
\[
\boxed{\lambda(\theta)=e\cos\theta+
\sqrt{R^2-e^2\sin^2\theta}}.
\]

\CorrigeQuestion{3}
La fonction précédente est évaluée sur un tour, pour
$0\leq\theta\leq2\pi$. La branche positive est retenue afin d'obtenir une
longueur de palette positive.
\begin{marginfigure}
\includegraphics[width=\linewidth]{10_02_c}
\end{marginfigure}

\CorrigeQuestion{4}
En dérivant par rapport au temps :
\[
\boxed{
\dot\lambda=
-e\dot\theta\sin\theta
-\frac{e^2\dot\theta\sin\theta\cos\theta}
{\sqrt{R^2-e^2\sin^2\theta}}
}.
\]
\begin{marginfigure}
\includegraphics[width=\linewidth]{10_03_c}
\end{marginfigure}

\CorrigeQuestion{5}
Pour une palette de section $S$, le débit algébrique instantané vaut
\[
\boxed{q(t)=S\,\dot\lambda(t)}.
\]
Le débit fourni est obtenu en ne conservant que la phase de refoulement.

\CorrigeQuestion{6}
Pour chaque valeur de l'excentration $e$, on calcule $\lambda(\theta)$, puis
$\dot\lambda(\theta)$ et enfin $q(\theta)=S\dot\lambda(\theta)$. Une
excentration plus importante augmente l'amplitude des variations de débit.
\begin{marginfigure}
\includegraphics[width=\linewidth]{10_04_c}
\end{marginfigure}

\CorrigeQuestion{7}
Pour cinq palettes régulièrement réparties, les contributions sont décalées
de $2\pi/5$. En ne conservant que les débits positifs, le débit total s'écrit
\[
\boxed{
q_{\mathrm{tot}}(\theta)=
\sum_{k=0}^{4}
\max\!\left(0,\,q\!\left(\theta-\frac{2k\pi}{5}\right)\right)
}.
\]
La somme des cinq contributions réduit fortement l'ondulation du débit.
\begin{marginfigure}
\includegraphics[width=\linewidth]{10_05_c}
\end{marginfigure}

\end{corrigebox}
